\documentclass[12pt,a4paper]{article}
\usepackage[utf8]{inputenc}
\usepackage[indonesian]{babel}
\usepackage{amsmath}
\usepackage{amsfonts}
\usepackage{amssymb}
\usepackage{geometry}
\usepackage{hyperref}

\geometry{margin=2.5cm}

\title{Metode Biseksi dalam Komputasi Numerik}
\author{Ahmad Fakhrul Bawani}
\date{}

\begin{document}

\maketitle

\section{Pendahuluan}
Metode Biseksi adalah salah satu algoritma pencarian akar paling sederhana untuk fungsi kontinu. Metode ini didasarkan pada Teorema Nilai Antara, yang menyatakan bahwa jika sebuah fungsi kontinu $f(x)$ memiliki tanda yang berbeda pada ujung-ujung interval $[a, b]$, maka setidaknya terdapat satu akar di dalam interval tersebut.

\section{Konsep dan Rumus}
Metode ini bekerja dengan cara membagi interval menjadi dua bagian secara berulang. Titik tengah $c$ dari interval $[a, b]$ dihitung dengan rumus:
\begin{equation}
c = \frac{a + b}{2}
\end{equation}

Setelah mendapatkan titik tengah, interval baru dipilih berdasarkan tanda dari nilai fungsi $f(c)$. Jika $f(a) \cdot f(c) < 0$, maka akar berada di sub-interval kiri $[a, c]$. Jika tidak, maka akar berada di sub-interval kanan $[c, b]$.

\section{Algoritma Iterasi}
\begin{enumerate}
  \item Tentukan interval $[a, b]$ sedemikian rupa sehingga $f(a) \cdot f(b) < 0$.
  \item Hitung titik tengah: $c = \frac{a + b}{2}$.
  \item Uji posisi akar:
  \begin{itemize}
    \item Jika $f(a) \cdot f(c) < 0$, maka akar di $[a, c]$, set $b = c$.
    \item Jika $f(b) \cdot f(c) < 0$, maka akar di $[c, b]$, set $a = c$
    \item Jika $f(c) = 0$, maka $c$ adalah akar eksak.
  \end{itemize}
  \item Ulangi langkah 2 dan 3 sampai lebar interval $|b - a|$ lebih kecil dari toleransi $\epsilon$ yang diinginkan.
\end{enumerate}

\section{Kelebihan dan Kekurangan}
\subsection{Kelebihan}
\begin{itemize}
  \item Selalu konvergen menuju akar (pasti ketemu selama fungsi kontinu).
  \item Sangat stabil dan tidak memerlukan perhitungan turunan fungsi.
  \item Mudah diimplementasikan dalam pemrograman.
\end{itemize}

\subsection{Kekurangan}
\begin{itemize}
  \item Konvergensi sangat lambat (memerlukan banyak iterasi untuk mencapai akurasi tinggi).
  \item Tidak dapat digunakan jika fungsi tidak berganti tanda di sekitar akar (misal pada akar ganda).
\end{itemize}

\end{document}